% GENERATED FILE. Edit canonical lessons, then run npm run generate:books.
% canonical-source-hash: fnv1a64:9ec1d9337f620a7c
% canonical-lessons: HI-C253-indradhanus, HI-C253-bhukamp, HI-C253-barh, HI-C253-purus, HI-C253-sisu

\chapter{A rainbow, An earthquake, A flood, A man, A baby}
\label{ch:hi-a-rainbow-an-earthquake-a-flood-a-man-a-baby}
\hlchaptermodality{253}

\begin{chapteropening}
\textbf{By the end of this chapter:} \emph{I can say a rainbow, an earthquake, a flood, a man and a baby.}

Complete the last lesson of chapter 253: I can say a rainbow, an earthquake, a flood, a man and a baby.
\end{chapteropening}

\section[\texorpdfstring{\hi{इंद्रधनुष} (indradhanuṣ)}{इंद्रधनुष (indradhanuṣ)}]{\hi{इंद्रधनुष} (indradhanuṣ) — a rainbow}

\label{lesson:HI-C253-indradhanus}



\begin{quote}
Before the new one: say the Hindi for a thorn, then the Hindi for mud.
\end{quote}

\subsection*{You'll want to know: \hi{इंद्रधनुष}}
\textbf{\hi{इंद्रधनुष}} — \emph{indradhanuṣ} — \textquotedblleft{}a rainbow\textquotedblright{}.

\begin{grammarlens}[title={What you've built}]
One more word for this chapter.
\end{grammarlens}

\subsection*{Your turn}
\begin{itemize}
\raggedright
  \item \emph{Say it:} \emph{indradhanuṣ}
  \item \emph{Say it:} \emph{indradhanuṣ}, once more
  \item \emph{Say it:} say \emph{indradhanuṣ}
  \item \emph{Recall:} say the Hindi for a thorn, then the Hindi for mud, then say \emph{indradhanuṣ} again
\end{itemize}

\begin{tcolorbox}[breakable,colback=teal!4,colframe=teal!35!black,title={Before you move on}]
What does \hi{इंद्रधनुष} mean? (\textquotedblleft{}A rainbow\textquotedblright{}.) Say it once more.
\end{tcolorbox}
\section[\texorpdfstring{\hi{भूकंप} (bhūkamp)}{भूकंप (bhūkamp)}]{\hi{भूकंप} (bhūkamp) — an earthquake}

\label{lesson:HI-C253-bhukamp}



\begin{quote}
Before the new one: say the Hindi for mud, then the Hindi for a rainbow.
\end{quote}

\subsection*{You'll want to know: \hi{भूकंप}}
\textbf{\hi{भूकंप}} — \emph{bhūkamp} — \textquotedblleft{}an earthquake\textquotedblright{}.

\begin{grammarlens}[title={What you've built}]
One more word for this chapter.
\end{grammarlens}

\subsection*{Your turn}
\begin{itemize}
\raggedright
  \item \emph{Say it:} \emph{bhūkamp}
  \item \emph{Say it:} \emph{bhūkamp}, once more
  \item \emph{Say it:} say \emph{bhūkamp}
  \item \emph{Recall:} say the Hindi for mud, then the Hindi for a rainbow, then say \emph{bhūkamp} again
\end{itemize}

\begin{tcolorbox}[breakable,colback=teal!4,colframe=teal!35!black,title={Before you move on}]
What does \hi{भूकंप} mean? (\textquotedblleft{}An earthquake\textquotedblright{}.) Say it once more.
\end{tcolorbox}
\section[\texorpdfstring{\hi{बाढ़} (bāṛh)}{बाढ़ (bāṛh)}]{\hi{बाढ़} (bāṛh) — a flood}

\label{lesson:HI-C253-barh}



\begin{quote}
Before the new one: say the Hindi for a rainbow, then the Hindi for an earthquake.
\end{quote}

\subsection*{You'll want to know: \hi{बाढ़}}
\textbf{\hi{बाढ़}} — \emph{bāṛh} — \textquotedblleft{}a flood\textquotedblright{}.

\begin{grammarlens}[title={What you've built}]
One more word for this chapter.
\end{grammarlens}

\subsection*{Your turn}
\begin{itemize}
\raggedright
  \item \emph{Say it:} \emph{bāṛh}
  \item \emph{Say it:} \emph{bāṛh}, once more
  \item \emph{Say it:} say \emph{bāṛh}
  \item \emph{Recall:} say the Hindi for a rainbow, then the Hindi for an earthquake, then say \emph{bāṛh} again
\end{itemize}

\begin{tcolorbox}[breakable,colback=teal!4,colframe=teal!35!black,title={Before you move on}]
What does \hi{बाढ़} mean? (\textquotedblleft{}A flood\textquotedblright{}.) Say it once more.
\end{tcolorbox}
\section[\texorpdfstring{\hi{पुरुष} (puruṣ)}{पुरुष (puruṣ)}]{\hi{पुरुष} (puruṣ) — a man}

\label{lesson:HI-C253-purus}



\begin{quote}
Before the new one: say the Hindi for an earthquake, then the Hindi for a flood.
\end{quote}

\subsection*{You'll want to know: \hi{पुरुष}}
\textbf{\hi{पुरुष}} — \emph{puruṣ} — \textquotedblleft{}a man\textquotedblright{}.

\begin{grammarlens}[title={What you've built}]
One more word for this chapter.
\end{grammarlens}

\subsection*{Your turn}
\begin{itemize}
\raggedright
  \item \emph{Say it:} \emph{puruṣ}
  \item \emph{Say it:} \emph{puruṣ}, once more
  \item \emph{Say it:} say \emph{puruṣ}
  \item \emph{Recall:} say the Hindi for an earthquake, then the Hindi for a flood, then say \emph{puruṣ} again
\end{itemize}

\begin{tcolorbox}[breakable,colback=teal!4,colframe=teal!35!black,title={Before you move on}]
What does \hi{पुरुष} mean? (\textquotedblleft{}A man\textquotedblright{}.) Say it once more.
\end{tcolorbox}
\section[\texorpdfstring{\hi{शिशु} (śiśu)}{शिशु (śiśu)}]{\hi{शिशु} (śiśu) — a baby, an infant}

\label{lesson:HI-C253-sisu}



\begin{quote}
Before the new one: say the Hindi for a flood, then the Hindi for a man.
\end{quote}

\subsection*{You'll want to know: \hi{शिशु}}
\textbf{\hi{शिशु}} — \emph{śiśu} — \textquotedblleft{}a baby, an infant\textquotedblright{}.

\begin{grammarlens}[title={What you've built}]
One more word for this chapter.
\end{grammarlens}

\subsection*{Your turn}
\begin{itemize}
\raggedright
  \item \emph{Say it:} \emph{śiśu}
  \item \emph{Say it:} \emph{śiśu}, once more
  \item \emph{Say it:} say \emph{śiśu}
  \item \emph{Recall:} say the Hindi for a flood, then the Hindi for a man, then say \emph{śiśu} again
\end{itemize}

\begin{tcolorbox}[breakable,colback=teal!4,colframe=teal!35!black,title={Before you move on}]
What does \hi{शिशु} mean? (\textquotedblleft{}A baby, an infant\textquotedblright{}.) Say it once more.
\end{tcolorbox}
